\honest{One Argument Against An Army}{Eliezer Yudkowsky}{2007}

Yesterday I wrote about reasoning in which no contrary argument is allowed. Here I suggest how people manage it: when they meet a contrary argument, they avoid lowering their confidence by rehearsing the support they already knew.

Suppose Freedonia suspects its neighbour Sylvania of causing a rash of meteor strikes. \nb{Freedonia, Sylvania and its ambassador Trentino come from the Marx Brothers film \textsc{Duck Soup} (1933).} The meteors struck near the border, the Sylvanian stock market moved before the strikes, and the ambassador muttered about ``heavenly vengeance.''

Someone points out that Sylvania trades with Freedonia for billions of dinars a year. You repeat your three points, and three outweigh one. Someone else points out that Sylvania has no space program. You repeat the three points again. Each time the balance is 3 to 1, and each time you feel more sure. After a hundred such objections you are surer still.

The problem is that you are counting the same evidence twice. Counting all the data twice would be bad enough, like a scientist who fails to get a significant result from 50 subjects and counts each subject twice. Counting only the favourable evidence twice is ``sheer farce,'' like the cartoon villain who divides the loot ``One for you, one-two for me.''

As I said yesterday, even a true belief must sometimes lose probability when new evidence comes in. Old support cannot offset the new evidence, because it was already counted.

``And yet it does appear to me'' that people facing a new counterargument search for a reason not to lower their confidence, and find the arguments they already knew. \nb{The post gives no evidence that this is common beyond the hypothetical Freedonia case.} I have to keep constant vigilance against doing it myself.

``With the right kind of wrong reasoning, a handful of support \ldots\ can stand off an army of contradictions.''
